\begin{figure}[htbp]
    \centering
    \begin{tikzpicture}[
        node distance=1.1cm and 2.0cm,
        box/.style={
            draw,
            rectangle,
            thick,
            minimum width=3.4cm,
            minimum height=1.1cm,
            align=center
        },
        smallbox/.style={
            draw,
            rectangle,
            thick,
            minimum width=3.1cm,
            minimum height=0.9cm,
            align=center
        },
        botharrow/.style={
            <->,
            >=Latex,
            thick
        }
    ]

        % Zentrale Komponente
        \node[box] (script) {Script};
        % Linke Eingaben
        % above left=1.0cm and 1.0cm
        \node[smallbox, left=1.2cm of script] (ci) {CI};
        \node[smallbox, below left=1.2cm and 1cmof script] (dev) {Entwickler\\(lokal)};

        % Rechte Komponenten
        \node[smallbox, right=2.2cm of script, yshift=1.4cm] (notes) {git notes/\\sonstiges};
        \node[smallbox, right=2.2cm of script, yshift=-0.4cm] (discovery) {testDiscovery.sh};
        \node[smallbox, below=0.5cm of discovery] (auditor) {testAuditor};
        \node[smallbox, below=0.5cm of auditor] (execution) {testExecution.sh};

        % Bidirektionale Verbindungen links
        \draw[botharrow] (ci.south) -- ++(0,-0.45) -| (script.north west);
        \draw[botharrow] (dev.east) -- (script.west);

        % Bidirektionale Verbindungen rechts
        \draw[botharrow] (script.north east) -- (notes.west);
        \draw[botharrow] (script.east) -- (discovery.west);
        \draw[botharrow] (script.south east) -- (auditor.west);
        \draw[botharrow] (script.south) -- (execution.west);
    \end{tikzpicture}
    \caption{Schematische Darstellung der Interaktion zwischen CI, lokalem Entwickler, Script und Hilfskomponenten}
    \label{fig:tool-overview}
\end{figure}